\textbf{Proof.}
Fix \(P\in\mathfrak P_1(N_0,K_T,K_D)\).  By \(d=r-1\) and \(r=1\), \(d=0\).  Equivalently, \(\mathcal C_D\) has rank zero by \(\mathrm{ass:exact-centered-rank}\).  The zero-dimensional convention in \(\mathrm{def:canonical-factorization}\) makes \(b_p(t)^\top z_m\) the empty sum, hence zero.  Substitution in \(\mathrm{ass:normalized-factor-model}\) gives, pointwise,
\[
 S^0_{p,m}(t)=s_p(t)+b_p(t)^\top z_m=s_p(t).
\]
Using \(\mathrm{def:target-curve}\) with \(p=1\) and \(m=\star\) therefore yields
\[
 \theta(t)=S^0_{1,\star}(t)=s_1(t).
\]

Now let \(\{P_\eta:|\eta|<\eta_0\}\) be a regular submodel of the stated kind.  Its target-pre marginal density, say \(q_{0,\star,\eta}\), is fixed in \(\eta\).  Therefore its target-pre score is
\[
 \left.\partial_\eta\log q_{0,\star,\eta}(O_{0,\star,j})\right|_{\eta=0}
 =\left.\partial_\eta\log q_{0,\star,0}(O_{0,\star,j})\right|_{\eta=0}=0.
\]
This is the precise ancillarity assertion; it does not assert that a path which also changes the target-pre marginal has zero target-pre score.

Finally fix \(r\ge2\) and a regular path \(z_{\star,\eta}=z_\star+\eta v\) which remains in the model.  Endpoint normalization is preserved because \(b_1(0)=0\) by \(\mathrm{def:canonical-factorization}\).  All donor coordinates are fixed, so the primitive donor operator and its exact rank and singular values are fixed.  Direct differentiation of \(\mathrm{def:target-curve}\) gives
\[
 \left.\partial_\eta\theta_\eta(t)\right|_{\eta=0}
 =\left.\partial_\eta\{s_1(t)+b_1(t)^\top(z_\star+\eta v)\}\right|_{\eta=0}
 =b_1(t)^\top v.
\]
The derivative is a nonzero element of \(C(\mathcal I)\) if and only if \(b_1^\top v\not\equiv0\), as claimed. \(\square\)